% =====================================================================
%  Capítulo 6 — Ver el sitio público y cerrar sesión
% =====================================================================
\chapter{Ver el sitio y cerrar sesión}\label{cap:salida}

\entradilla{El final del recorrido: comprobar lo que ve el cliente y salir del panel de forma
segura.}

\section{Ver el sitio público}\label{sec:sitio-publico}

Al final de la barra lateral, \textbf{\icExterno\ Ver sitio público} abre la página de inicio de Training~InnovaQuote en una pestaña nueva. El panel sigue abierto en la pestaña original.

Úselo para comprobar sus cambios desde el punto de vista del cliente. Por ejemplo:

\begin{itemize}
  \item Que una unidad de negocio nueva aparezca en \msg{Elija el tipo de servicio}, después de
        crear sus preguntas.
  \item Que el formulario de una unidad muestre las preguntas y los bloques en el orden
        esperado.
\end{itemize}

\captura[8cm]{06-01-sitio-publico}{/}{Página de inicio de Training~InnovaQuote.}
  {la página de inicio pública de Training~InnovaQuote, abierta desde «Ver sitio público».}

\begin{importante}
Si envía una solicitud de prueba desde el sitio público, se creará una cotización real con su
código y llegarán los correos de confirmación. Para revisar el formulario sin crear solicitudes,
use \botonsec{\icPlay\ Probar como cliente} (\seccion{sec:probar}).
\end{importante}

\section{Cerrar sesión}\label{sec:cerrar-sesion}

\begin{pasos}
  \item Guarde lo que esté editando. Por ejemplo, en el árbol de decisión, termine con
        \boton{Guardar pregunta} o \botonsec{Cancelar}.
  \item Vaya al final de la barra lateral. En un celular, abra antes el menú con
        \textbf{\icMenu} (\seccion{sec:movil}).
  \item Debajo de \textbf{Ver sitio público} verá \textbf{\icSalir\ Cerrar sesión}, con el
        correo de su cuenta.
  \item Haga clic en \textbf{Cerrar sesión}.
\end{pasos}

\capturaconnotas[6.5cm]{06-02-cerrar-sesion}{/dashboard}{Opción «Cerrar sesión» al final de la barra lateral.}
  {un recorte de la parte inferior de la barra lateral, con «Ver sitio público», «Cerrar
   sesión» y el correo de la cuenta.}
  {\textbf{Antes de salir}

   \begin{itemize}
     \item Compruebe que el último aviso fue de éxito.
     \item Si dejó una pregunta abierta en el árbol de decisión, guárdela o cancélela: al cerrar
           sesión, el sistema no le avisa y los cambios sin guardar se pierden.
     \item El correo bajo \textbf{Cerrar sesión} le confirma con qué cuenta está trabajando.
   \end{itemize}}

\begin{resultado}
La sesión se cierra y el navegador lo lleva a la página de inicio pública de Training~InnovaQuote. Para
volver a entrar, siga los pasos de la \seccion{sec:iniciar-sesion}.
\end{resultado}

\begin{consejo}
\begin{itemize}
  \item Cierre sesión siempre que use un computador compartido. Cerrar la pestaña o el
        navegador no basta, sobre todo si marcó \campo{Mantener la sesión iniciada}.
  \item Si deja el panel sin usar por un tiempo largo (por defecto, 2 horas), la sesión se
        cierra sola y deberá iniciar sesión de nuevo. Si estaba llenando un formulario, podría
        perder lo que escribió.
\end{itemize}
\end{consejo}

\newcommand{\caprefclaro}[1]{\hyperref[#1]{\textcolor{cielo}{\fontseries{sb}\selectfont capítulo~\ref*{#1}}}}

\begin{tcolorbox}[enhanced, colback=noche, colframe=noche, coltext=white, arc=6pt,
    left=16pt, right=16pt, top=14pt, bottom=14pt, before skip=18pt,
    overlay={\begin{scope}
        \clip (frame.south west) rectangle (frame.north east);
        \pic[opacity=0.35] at ([xshift=-1.6cm, yshift=-0.3cm]frame.north east) {anillo={2.6cm/1.2pt/cielo}};
      \end{scope}}]
  {\montserrat\bfseries\large\color{white}Recorrido completo}\par\vspace{4pt}
  {\color{niebla}Ya conoce el flujo completo del administrador:}\par\vspace{6pt}
  \begin{itemize}[label={\color{cielo}\ico{F00C}}, leftmargin=1.6em]
    \item Iniciar sesión y recuperar la contraseña (\caprefclaro{cap:acceso}).
    \item Revisar solicitudes, enviar propuestas y consultar clientes (\caprefclaro{cap:operacion}).
    \item Mantener el catálogo y el árbol de decisión (\caprefclaro{cap:catalogo}).
    \item Administrar usuarios, invitaciones y configuración (\caprefclaro{cap:sistema}).
    \item Revisar el sitio público y cerrar sesión (\caprefclaro{cap:salida}).
  \end{itemize}
\end{tcolorbox}
